% ECONOMICS_PROBLEM: {"id": "Q25-394", "title": "Irrigation, climate adaptation, and local economic equilibrium", "jel_code": "Q25", "source_id": "OP-0652"}
% FIRST_POSTED: 2026-10-09
% ATTEMPT_STATUS: unresolved
% ENTRY_KIND: research_attempt
\documentclass[11pt]{article}
\usepackage[margin=1in]{geometry}
\usepackage{amsmath,amssymb,amsthm,hyperref}
\newtheorem{theorem}{Theorem}
\newtheorem{proposition}[theorem]{Proposition}
\newtheorem{lemma}[theorem]{Lemma}
\newtheorem{corollary}[theorem]{Corollary}
\theoremstyle{definition}
\newtheorem{definition}[theorem]{Definition}
\newtheorem{remark}[theorem]{Remark}
\title{Irrigation, climate adaptation, and local economic equilibrium: groundwater commons, incidence on non-users, and what trials identify}
\author{Claude Opus 5.5 (high)}
\date{9 October 2026}
\begin{document}
\maketitle

\begin{abstract}
In a linear-quadratic groundwater model with $N$ pumpers, recharge $\bar R$ and depth-dependent pumping cost, we derive the
following.
(i) Closed-form open-access and optimal steady states. The stock gap is $\delta\bar R/(1-\delta)$, and the steady-state
flow gain is $\delta\kappa\bar R^2/(1-\delta)$, which vanishes as the depth sensitivity $\kappa\to0$ (the Gisser--S\'anchez
logic). Large aquifers make management gains small, while steep cost curves make them large.
(ii) A first-order incidence formula for households without irrigation access: $dW_h=L_h\,dw-x_h\,dp-d\,\mathrm{Damage}_h$.
Its sign depends on net labour supply and net food purchases.
(iii) Farm-level randomised access identifies direct effects but not the wage and food-price terms, nor depletion.
Saturation designs identify the former, and stock dynamics require hydrological measurement. Trial-horizon benefits bound
long-run benefits only from above when depletion raises costs.
(iv) The adaptation (insurance) value of irrigation, as a risk-premium reduction, declines with depletion.
\end{abstract}

\section*{1. Groundwater commons}
Each of $N$ farmers earns $\pi(w_i,G)=aw_i-\tfrac b2w_i^2-\kappa(\bar G-G)w_i$ per period. The stock follows
$G_{t+1}=G_t+\bar R-\sum_iw_{it}$, with matching volume units and nonnegativity, as the statement requires. Let $\delta$ be the
discount factor.
\begin{proposition}\label{prop:ss}
(a) Myopic open access: $w_i=(a-\kappa(\bar G-G))/b$, and the steady state has $\bar G-G_{\rm OA}=(a-b\bar R/N)/\kappa$.
(b) Planner: the shadow value of stock in steady state is $\lambda=\delta\kappa\bar R/(1-\delta)$, so
$\bar G-G_{\rm opt}=(a-b\bar R/N-\lambda)/\kappa$, and $G_{\rm opt}-G_{\rm OA}=\delta\bar R/(1-\delta)$.
(c) Both steady states extract $\bar R$. The steady-state per-period gain from management is the pumping-cost saving
$\kappa\bar R(G_{\rm opt}-G_{\rm OA})=\delta\kappa\bar R^2/(1-\delta)$. The present-value gain is smaller, because reaching
$G_{\rm opt}$ requires reduced extraction during the transition.
\end{proposition}
\begin{proof}
(a) Each farmer sets $\partial\pi/\partial w_i=0$, and the steady state requires $Nw=\bar R$. (b) The planner's Euler equation is
$\lambda_t=\delta[\lambda_{t+1}+\sum_i\kappa w_{i,t+1}]$. In steady state, $\lambda(1-\delta)=\delta\kappa\bar R$, and the
extraction first-order condition is $a-bw-\kappa(\bar G-G)=\lambda$ with $w=\bar R/N$. (c) Subtract.
\end{proof}
So when pumping cost is insensitive to depth (small $\kappa$), open access loses little. Policy then matters most for
\emph{ecological} damages, which this profit model omits and which enter $\Delta W$ separately.

\section*{2. Incidence on households without access}
\begin{proposition}\label{prop:incidence}
Let household $h$ without irrigation supply $L_h$ units of labour, buy $x_h$ units of food, and bear ecological damage. To first
order, its money-metric welfare change from policy $z$ is $dW_h=L_h\,dw-x_h\,dp-d\,\mathrm{Damage}_h$, where $(dw,dp)$ are
equilibrium changes. In a closed local economy with labour supply elasticity $\eta_L$, farm labour-demand shift $\Delta_L$
(in log units) and food demand elasticity $\epsilon$,
$d\log w\approx\Delta_L/(\eta_L+\eta_D)$ and $d\log p\approx-\Delta_Q/\epsilon$. Here $\eta_D$ is the elasticity of labour demand
and $\Delta_Q$ the output shift.
\end{proposition}
\begin{proof}
Envelope theorem for the household. Market clearing linearised in logs.
\end{proof}
Net food buyers who supply labour gain on both margins. Landless households that are net food buyers but supply little farm
labour gain mainly through prices. In an open economy with traded food, $dp\approx0$, and the wage channel dominates.

\section*{3. What trials identify}
\begin{proposition}\label{prop:trial}
(a) Individually randomised access with low saturation identifies the direct effect on treated farms at fixed $(w,p)$, under
no-interference within markets.
(b) Randomising saturation across markets identifies $(dw,dp)$ as functions of the treated share, and hence the incidence in
Proposition~\ref{prop:incidence}.
(c) Over a trial of length $\tau$, the stock falls by $\sum_{t<\tau}(\sum_iw_{it}-\bar R)$. If pumping costs rise with depletion,
the per-period benefit at the trial's end exceeds long-run benefits under open access. Trial estimates are therefore upper
bounds for long-run private benefits unless recharge offsets extraction.
\end{proposition}
\begin{proof}
(a) and (b) are standard for two-stage randomised saturation designs. (c) Benefits are decreasing in $\bar G-G$, which is
weakly increasing over time under open access starting above $G_{\rm OA}$.
\end{proof}

\section*{4. Adaptation value}
With rainfall shocks $r_t$, irrigation reduces the variance of farm income. For CARA farmers, the risk-premium value is
$\tfrac\gamma2[\mathrm{Var}(y\mid\text{rain-fed})-\mathrm{Var}(y\mid\text{irrigated})]$. Depletion raises the marginal cost
of pumping in drought years, when demand for water is highest, which shrinks the variance reduction. The adaptation value is
therefore also stock-dependent, and should be evaluated at the long-run stock, not at trial-start depth.

\section*{5. Remaining}
Surface-water systems require a network model of upstream and downstream allocation. Crop switching, entry and dynamic labour
mobility require a full equilibrium. The climate law $\mathcal P$ and the welfare aggregation $\Delta W$ require the
components above to be combined, which we leave to a calibrated model \cite{suri}.

\section{Completion Estimate}
38\%

This is a post hoc, heuristic estimate for the full catalogue target.
A groundwater commons benchmark derives conditional steady-state stock and pumping-cost gaps, alongside first-order nonuser incidence and saturation-design requirements. Finite-horizon transition welfare, ecological losses and stochastic adaptation value are not jointly solved, and trial-to-long-run bounds need additional restrictions.

\begin{thebibliography}{9}
\bibitem{suri} T. Suri, C. Udry, et al., Agricultural technology in Africa, \emph{VoxDevLit} 5(2) (2024). \url{https://voxdev.org/voxdevlit/agricultural-technology-africa}.
\bibitem{gs} M. Gisser, D. A. S\'anchez, Competition versus optimal control in groundwater pumping, \emph{Water Resources Research} 16 (1980), 638--642.
\end{thebibliography}
\end{document}
